\section{保证合同}
\begin{enumerate}
  \item 绝对的难点。
  \item \zd{单务、无偿}。
  \item 公益为目的的非营利法人、组织不得做保证人（承担无限责任）\yjt 保证合同无效，承担缔约过失责任。
  \item 订立方式。\zd{表明自己要承担责任}。
        \begin{enumerate}
          \item 单独签一份书面保证合同。
          \item 主合同中有保证条款，第三人签字。
          \item 无条款，第三人以保证人签字。
          \item 单方出示\zd{保证/担保函}。考试就考这种。
        \end{enumerate}
  \item \zd{保证的方式}。\zd{保证的标的物是保证人的信誉}（人保）\yjt 只有人保有承担责任的顺序问题。\al{如甲到期不向乙还款，由本人以某处私房抵债 ×}。
        \begin{enumerate}
          \item 有顺序 \yjt \zd{一般保证}（核心考点）\yjt \zd{先诉抗辩权}（独有的，必考）。
                \begin{enumerate}
                  \item 程序上不能单独先起诉。
                  \item 实体上不得单独先执行。
                \end{enumerate}
          \item 没顺序 \yjt \zd{连带责任保证}。\zd{必须明示，不能推定}。
        \end{enumerate}
  \item 保证范围。有约从约 \yjt 全部（包括利息等）。
  \item 保证的从属性 = 担保从属性。减轻/加重保证责任 \yjt 通知/书面同意。
        \begin{enumerate}
          \item \zd{不允许砍头息}。先在本金中扣除利息，以实际数额为准。\al{借 11W，约定其中 1W 是利息，实际数额 10W}。
          \item 履行期限。\zd{延长期限覆盖全部保证期间的，保证人无须承担保证责任}。\al{2011/1/1 到期，保证期 1 年，约定借款延期到 2012/1/1}。
          \item 债权转让。\zd{内外有别，双重通知}。例外。有约从约（\al{如约定禁止转让}）。
          \item 主债务转让。
                \begin{enumerate}
                  \item 加入 \yjt 通知。
                  \item 退出 \yjt \zd{双重同意且有顺序}（从随主）。主债权人（不要式）\yjt 保证人（书面同意）。
                  \item 不利后果。保证人不再承担相应的保证责任（转多少免多少）。
                  \item \al{该债务转让应取得担保人同意 ×}。担保人的同意指的是同意了一起承担，不同意不承担。而不是指对\zd{债务转让协议效力本身}的任何肯定 \hz 否定。
                \end{enumerate}
          \item 保证期间。保证人承担责任的期间。
                \begin{enumerate}
                  \item 不变期间 \yjt 因为无限责任。\zd{考试采用除斥期间说}。
                  \item 有约从约，\zd{无约 6 个月}，约定不明按无约。\al{直至本息还清 \yjt 6 个月}。
                  \item 多段保证，分段独立计算。\al{7/3 还 10W，8/3 还 20W，9/3 还 30W，甲对 3 笔还款均提供保证}。
                  \item \zd{一旦债权人起诉，保证期间消灭，保证人拒绝权消灭 \yjt 主张保证责任 3 年诉讼时间}。
                \end{enumerate}
        \end{enumerate}
        % q=WzAsMTMsWzAsMCwi5YqhXFxcXOWtnyJdLFs0LDAsIuadg1xcXFzpqawiXSxbMiwxLCLmm7lcXFxc5L+dIl0sWzAsMiwi5Li777ya5YCf5qy+5ZCI5ZCMIl0sWzEsMiwiMjAyMi4xLjEiXSxbMiwyLCIyMDIzLjEuMSIsWzI0MCw2MCw2MCwxXV0sWzUsMiwiMjAyNi4xLjEiXSxbMCwzLCLku47vvJrkv53or4HlkIjlkIwiXSxbMSwzLCLkv53or4HmnJ/pl7QiXSxbMiwzLCIyMDIzLjEuMiIsWzI0MCw2MCw2MCwxXV0sWzUsMywiMjAyMy43LjEiLFsyNDAsNjAsNjAsMV1dLFswLDEsIuWPr+WPmOacn+mXtCIsWzAsNjAsNjAsMV1dLFswLDQsIuS4jeWPmOacn+mXtCIsWzAsNjAsNjAsMV1dLFswLDEsIuS4u1xcXFwyMDAwV1xcXFzlgJ/mrL7lkIjlkIwiLDFdLFsyLDEsIuS/neivgeWQiOWQjFxcXFzku44iLDNdLFszLDRdLFs0LDVdLFs1LDYsIuivieiuvOaXtuaViOacn+mXtFxcXFwzIOW5tCIsMV0sWzcsOF0sWzgsOV0sWzksMTAsIuS/neivgeacn+mXtFxcXFw2IOS4quaciCIsMSx7fSxbMjQwLDYwLDYwLDFdXSxbNywxMl0sWzMsMTFdXQ==
\[
  \begin{tikzcd}[row sep=small]
    & \txt{务\\孟} &&& \txt{权\\马} \\
    & \textcolor{rgb,255:red,214;green,92;blue,92}{\txt{可变期间}} & \txt{曹\\保} \\
    & \txt{主。借款合同} & \txt{2022.1.1} & \txt{2023.1.1\\借期届满} &&& \txt{2026.1.1} \\
    \textcolor{rgb,255:red,214;green,92;blue,92}{\txt{不变期间}} & \txt{从。保证合同} & \txt{保证期间} & \txt{2023.1.2} &&& \txt{2023.7.1} \\
    & \txt{从。保证合同} & \txt{债权人行权} & \txt{2023.5.2\\拒绝权消灭} &&& \txt{2024.5.1}
    \arrow["\txt{主\\2000W\\借款合同}"{description}, from=1-2, to=1-5]
    \arrow["\txt{保证合同\\从}"{marking}, from=2-3, to=1-5]
    \arrow[from=3-2, to=3-3]
    \arrow[from=3-3, to=3-4]
    \arrow[from=4-2, to=4-3]
    \arrow[from=4-3, to=4-4]
    \arrow[from=4-2, to=4-1]
    \arrow[from=3-2, to=2-2]
    \arrow[from=5-2, to=5-3]
    \arrow[from=5-3, to=5-4]
    \arrow["\txt{诉讼时效期间}", from=3-4, to=3-7]
    \arrow["\txt{3 年}"', from=3-4, to=3-7]
    \arrow["\txt{保证期间}", from=4-4, to=4-7]
    \arrow["\txt{3 年}"', from=4-4, to=4-7]
    \arrow["\txt{保证合同诉讼时效}", from=5-4, to=5-7]
    \arrow["\txt{3 年}"', from=5-4, to=5-7]
  \end{tikzcd}
\]

  \item 保证人权利。
        \begin{enumerate}
          \item 追偿权。<= 主债权的范围。\al{甲向乙采购 10 袋菊花茶，货款 10W，丙提供担保。有 2 袋无法饮用，价值 2 W，丙付款 10W}。
          \item \zd{法定代位权}。优先于债务人的其他债权人。追偿权的附带品。
          \item 抵销权、撤销权。债务人有的他都有。考的很少。
          \item 放弃诉讼时效抗辩/先诉抗辩权，丧失（从属性）/不丧失对债务人的追偿权。
        \end{enumerate}
  \item 主债务人。最终责任人。
  \item 共同担保。保证人 >= 2 个。金字塔尖的难点，必考。完全符合基本原理，没有例外。
        \begin{enumerate}
          \item 按标的物。共同人保、共同物保、混合担保。
          \item 考试考起诉时（是否有一般担保）、执行时（是否有债务人的物保 \yjt 应当先）、追偿时（\zd{是否连带共同}）怎么操作。
          \item 范围。有约从约，无约是第 3 种情况。
                \begin{enumerate}
                  \item 外部一般保证，内部按份。
                  \item 外部一般保证，内部连带。默认。
                  \item 外部连带保证，内部按份。
                  \item 内外均连带/双重连带。
                \end{enumerate}
          \item 清偿顺序。
                \begin{enumerate}
                  \item 有约从约 \yjt \yjt 有债务人的物保 \zd{应当先}就债务人的物保实现债权（先找保证人不公平）。
                  \item \zd{主张全部债权的}。有约从约 \yjt 所有按比例而不是\al{应先处理 xx，再处理 xx，是一起全部处理}。
                \end{enumerate}
          \item 内部追偿问题。\zd{原则不让互相追，尽可能找主债务人}。
                \begin{enumerate}
                  \item 有约从约。但有内无外。
                  \item 互相不知道对方存在。无追偿权。\al{各自、分别}。
                  \item 知道。
                        \begin{enumerate}
                          \item 按份共同。无追偿权。自己承担自己的份额。
                          \item 连带共同。先债务人 \yjt 不足部分按比例找其他担保人分担 \yjt 债务人。
                                \al{甲债务人，乙丙丁担保。乙 100W 追偿，甲 10W，剩余 90W，向丙丁各自追 30W。乙丙丁各自的 30W 继续找甲}。
                        \end{enumerate}
                \end{enumerate}
        \end{enumerate}
  \item \zd{追偿权有风险} \yjt \zd{反担保} \yjt 担保的是担保人的追偿权。债务人 \hz 第三人提供。
\end{enumerate}
